\documentclass[10pt,a4paper]{article}
\usepackage[T1]{fontenc}
\usepackage[utf8]{inputenc}
\usepackage{lmodern}
\usepackage{amsmath,amssymb}
\usepackage[margin=21mm]{geometry}
\usepackage{microtype}
\usepackage[hidelinks]{hyperref}
\usepackage{enumitem}
\setlist[itemize]{leftmargin=*,nosep,topsep=3pt}
\setlength{\parindent}{0pt}
\setlength{\parskip}{5pt}
\setlength{\emergencystretch}{3em}
\newcommand{\question}[1]{\par\medskip{\normalsize\bfseries #1}\par\smallskip}
\pagestyle{plain}

\begin{document}
{\Large\bfseries Neurophysiological mechanisms of deep brain stimulation across spatiotemporal resolutions\par}

\textbf{Paper:} Neumann, Steiner and Milosevic (2023).\quad Review and physiological synthesis.\par
\href{https://doi.org/10.1093/brain/awad239}{DOI: \nolinkurl{10.1093/brain/awad239}}\par
{\footnotesize Page references below use PDF page numbers, including front matter.}

\question{To which general question does this paper relate?}
What DBS acts on and what can be measured during stimulation. Useful for selecting a defensible actuator model and avoiding a universal excitation/inhibition assumption.

\question{How does it answer the general questions?}
\textbf{Mechanisms.} Electrical stimulation recruits axons and propagates activity through circuits; local firing can nevertheless be suppressed. Synaptic depression and the balance of excitatory/inhibitory inputs depend on the target and stimulation frequency. Thus a local firing-rate effect does not characterize the entire network response (PDF pp.~2--7).

\textbf{Measurements and objective.} Unit activity, spontaneous LFP rhythms and evoked responses provide complementary observations. Beta suppression accompanies motor improvement in several cited studies, but its precise mechanistic role remains disputed. Evoked responses may also inform programming (PDF pp.~6--9).

\textbf{Model relevance.} A constant additive input is a simplification requiring validation. Frequency-dependent synaptic effects may require extra dynamic states or uncertain gains, rather than only a memoryless monotone nonlinearity. The review provides no unique numerical delayed-model parameterization and does not establish beta suppression as a sufficient clinical objective.

\question{What papers does it cite that could provide more information?}
\begin{itemize}
\item Yu et al. (2020), \emph{Frequency-Specific optogenetic deep brain stimulation of subthalamic nucleus improves parkinsonian motor behaviors}: intervention evidence (ref.~44, PDF p.~11).
\item Steiner et al. (2019), \emph{Connectivity and dynamics underlying synaptic control of the subthalamic nucleus}: synaptic mechanisms (ref.~60, PDF p.~11).
\end{itemize}

\end{document}
